\chapter*{\centering \fontsize{14pt}{15pt}{\khmerfontmoul អត្ថបទសង្ខេប}}
\addcontentsline{toc}{chapter}{\fontsize{12pt}{15pt}{\khmerfontmoul អត្ថបទសង្ខេប}}

% Khmer body: Khmer OS Battambang with widened leading (stacked Khmer glyphs
% need ~1.6x); emergencystretch gives the justifier room because Khmer breaks
% only at the phrase spaces present in the text. Latin/math runs inherit the
% font cleanly (Battambang ships Latin glyphs; math stays in the math font).
\begingroup
\khmerfont\fontsize{12pt}{19pt}\selectfont
\setlength{\emergencystretch}{3em}

\hs ការវាយតម្លៃ និងទទួលធានា (underwriting) ផ្នែកធានារ៉ាប់រងសុខភាពនៅក្នុងទីផ្សារ កំពុងលេចឡើង (emerging markets) ដូចជាប្រទេសកម្ពុជា កំពុងប្រឈមនឹងបញ្ហាប្រឈម រចនាសម្ព័ន្ធ រួមមាន អត្រាជ្រៀតចូលនៃធានារ៉ាប់រងទាប ហេដ្ឋារចនាសម្ព័ន្ធថែទាំសុខភាព ដែលនៅខ្ចាត់ខ្ចាយ និងកំណត់ត្រាសុខភាពតាមរយៈពេលវែង (longitudinal) ដែលមានកម្រិត។ ប្រព័ន្ធ underwriting ផ្អែកលើច្បាប់ឋិតិវន្ត (static rule-based) បែបប្រពៃណី គឺមិនមាន ភាពល្អប្រសើរបំផុត ដោយសារវាមិនគិតពិចារណាអំពីអន្តរកម្មរវាងលក្ខណៈ (feature interactions) មិនអាចសម្របខ្លួនទៅនឹងការប្រែប្រួលនៃផតហ្វូលីយ៉ូ (portfolio drift) និងអាចបង្ករការរើសអើងដោយអចេតនាចំពោះក្រុមប្រជាសាស្ត្រមួយចំនួន។ និក្ខេបបទនេះ ស្រាវជ្រាវថាតើ contextual bandits --- ដែលជាប្រភេទនៃក្បួនដោះស្រាយរៀនតាមអនឡាញ (online learning) សម្រាប់ការសម្រេចចិត្តតាមលំដាប់ក្រោមភាពមិនច្បាស់លាស់ --- អាច ជំនួសច្បាប់ឋិតិវន្តនៅក្នុងការវាយតម្លៃ underwriting ផ្នែកធានារ៉ាប់រងសុខភាពនៅទីផ្សារ កំពុងលេចឡើងបានឬអត់។

ការស្រាវជ្រាវនេះរចនា និងអនុវត្តក្របខ័ណ្ឌ contextual bandit ដោយប្រើក្បួនដោះស្រាយ ចំនួនបី (LinUCB, LinTS និង Epsilon-Greedy) ហើយប្រៀបធៀបពួកវាជាមួយនឹង គោលជំហរមូលដ្ឋាន (baseline) នៃច្បាប់ឋិតិវន្ត XGBoost លើសំណុំទិន្នន័យសំយោគ (synthetic) នៃអ្នកដាក់ពាក្យសុំធានារ៉ាប់រងសុខភាពកម្ពុជាចំនួន 2{,}000 នាក់ ដែលផ្អែក លើការស្ទង់មតិប្រជាសាស្ត្រ និងសុខភាពកម្ពុជា ឆ្នាំ 2021--22 \citep{NIS2023}។ ឧបករណ៍ក្លែងធ្វើរង្វាន់តាមបែបអាក់ទ្យូអារី (actuarial) ផ្អែកលើប្រាក់ចំណេញ វាយតម្លៃ សកម្មភាព underwriting ចំនួនបួន (ស្តង់ដារ, ដាក់អត្រា (rated), បដិសេធ (decline), និងបញ្ជូនបន្ត (refer))។ របាំងការពារ (guardrails) បែប sliding-window នៃសន្ទស្សន៍ ស្ថិរភាពចំនួនប្រជាជន (Population Stability Index, PSI) និងច្បាប់បួនភាគប្រាំ (four-fifths rule) របស់ EEOC សហរដ្ឋអាមេរិក ត្រួតពិនិត្យរួមគ្នានូវភាពយុត្តិធម៌តាម តំបន់ និងតាមមុខរបរ។ លទ្ធផលសំខាន់ៗទាំងអស់ត្រូវបានរាយការណ៍ជាមធ្យមភាគលើ seeds ឯករាជ្យចំនួន 20 ជាមួយនឹងតេស្ត Wilcoxon signed-rank បែបផ្គូផ្គង (paired) និងចន្លោះទំនុកចិត្ត $95\%$ បែប bootstrap; EXP-008 (human-in-the-loop, §5.4) ត្រូវបានរាយការណ៍លើ seeds ឯករាជ្យចំនួន 20 (Wilcoxon $p < 0.001$)។

ការពិសោធន៍ដែលមានការគ្រប់គ្រងចំនួនបួន បានធ្វើសុពលកម្មក្របខ័ណ្ឌនេះ។ EXP-005 បង្ហាញថា LinUCB សម្រេចបានរង្វាន់បង្គរ (cumulative reward) ចំនួន \$90{,}540 ក្នុងរយៈពេល 5{,}000 ជុំ ធៀបនឹង \$72{,}292 របស់គោលជំហរមូលដ្ឋានឋិតិវន្ត ($+25\%$, Wilcoxon $p < 0.001$, Cohen's $d = 2.98$) ហើយកាត់បន្ថយ regret ជាមធ្យមនៅក្នុង 500 ជុំចុងក្រោយ ពី \$9.01 មកត្រឹម \$2.20 ក្នុងមួយជុំ។ EXP-006 បញ្ជាក់ថា PSI បែប sliding-window អតិបរមា នៅតែស្ថិតក្នុងតំបន់ GREEN/AMBER សម្រាប់ទាំងតំបន់ (0.082) និងមុខរបរ (0.123) ខណៈដែលច្បាប់បួនភាគប្រាំរបស់ EEOC ត្រូវបានបំពេញ ដោយមានសមាមាត្រភាពស្មើគ្នា (parity ratios) $85.7\%$ និង $90.1\%$ រៀងៗខ្លួន។ EXP-007 ដាក់ចំណាត់ថ្នាក់ក្បួនដោះស្រាយតាម regret បង្គរ ដែល LinTS និង LinUCB មិនអាចបែងចែកគ្នាបានតាមស្ថិតិ ហើយទាំងពីរសុទ្ធតែយកឈ្នះយ៉ាងច្បាស់ លាស់លើគោលជំហរមូលដ្ឋានឋិតិវន្ត និងគោលជំហរស្វែងរកឯកសណ្ឋាន (uniform-exploration)។ EXP-008 រុំ LinUCB ក្នុងស្រទាប់ human-in-the-loop ហើយបង្កើនរង្វាន់បង្គរ $14.8\%$ លើគោលជំហរមូលដ្ឋាន mathematical-REFER ដោយចំណាយលើការត្រួតពិនិត្យ ដោយមនុស្សត្រឹម $2.2\%$ នៃរង្វាន់។ ការវិភាគភាពរឹងមាំ (robustness) ក្រោយការ ពិសោធន៍ (post-hoc) បានកំណត់ពិដានសមត្ថភាពតាមទ្រឹស្តីមួយ៖ គោលនយោបាយថេរ ដែលមិនអាចទទួលយកបាន (inadmissible constant policy) ឈ្មោះ AlwaysRATED ដែលយកឈ្នះ bandits ទាំងពីរដោយ 30--33\% --- ដែលកំណត់ព្រំដែន ប៉ុន្តែមិន ធ្វើឱ្យលុបចោលលទ្ធផលនេះឡើយ។

លទ្ធផលទាំងនេះផ្តល់នូវភស្តុតាងបញ្ជាក់គំនិត (proof of concept) ដ៏ម៉ឺងម៉ាត់ថា contextual bandits អាចយកឈ្នះការវាយតម្លៃ underwriting ផ្អែកលើច្បាប់ឋិតិវន្ត ដែលវានឹងជំនួស ទាំងផ្នែកប្រាក់ចំណេញ ខណៈពេលដែលនៅតែរក្សាបាននូវភាពយុត្តិធម៌ ខាងប្រជាសាស្ត្រ --- ហើយនៅតែមានភាពស្រាលក្នុងការគណនា ដោយសារគំរូមូលដ្ឋានគឺ ជាឧបករណ៍ប៉ាន់ស្មានលីនេអ៊ែរ (linear estimator) ដែលមានលក្ខណៈ 34 និងត្រូវបាន ធ្វើបច្ចុប្បន្នភាពក្នុងទម្រង់បិទ (closed form); ពិដានថេរដែលមិនអាចទទួលយកបានដែល បានកំណត់ខាងលើ កំណត់ព្រំដែន ប៉ុន្តែមិនលុបចោលនូវការរួមចំណែកនេះធៀបនឹងជម្រើស ដែលអាចទទួលយកបាន (admissible alternatives) នោះឡើយ។ និក្ខេបបទនេះរួមចំណែក នូវ harness ពិសោធន៍ 20-seed ដែលអាចផលិតឡើងវិញបាន (reproducible) សំណុំ ទិន្នន័យសំយោគដែលបានក្រិតតាមបរិបទកម្ពុជា ការត្រួតពិនិត្យភាពយុត្តិធម៌ដោយ PSI បែប sliding-window និងច្បាប់បួនភាគប្រាំរបស់ EEOC ព្រមទាំង wrapper បែប human-in-the-loop ដែលអាចផ្តល់ព័ត៌មានដល់ការស្រាវជ្រាវនាពេលអនាគត និងការ អនុវត្តក្នុងឧស្សាហកម្មនៃធានារ៉ាប់រងផ្អែកលើក្បួនដោះស្រាយ (algorithmic insurance) នៅទីផ្សារកំពុងលេចឡើង។
\endgroup
